\documentclass[12pt,a4paper]{article}

\usepackage[utf8]{inputenc}
\usepackage[T1]{fontenc}
\usepackage[brazil]{babel}
\usepackage[top=3cm,left=3cm,bottom=2cm,right=2cm]{geometry}
\usepackage{mathptmx}
\usepackage{inconsolata}
\usepackage{setspace}
\usepackage{indentfirst}
\usepackage{titlesec}
\usepackage{booktabs}
\usepackage{array}
\usepackage{listings}
\usepackage{xcolor}
\usepackage{xurl}
\urlstyle{same}
\usepackage{tikz}
\usetikzlibrary{arrows.meta}

\onehalfspacing
\setlength{\parindent}{1.25cm}

\titleformat*{\section}{\normalfont\bfseries\normalsize}
\titleformat*{\subsection}{\normalfont\bfseries\normalsize}

\lstset{
    basicstyle=\ttfamily\footnotesize,
    breaklines=true,
    columns=fullflexible,
    keepspaces=true,
    frame=single,
    rulecolor=\color{black!25},
    backgroundcolor=\color{black!4},
    xleftmargin=0.2cm,
    xrightmargin=0.2cm,
    framexleftmargin=4pt,
    framexrightmargin=4pt,
    framextopmargin=4pt,
    framexbottommargin=4pt,
    literate=
        {á}{{\'a}}1 {é}{{\'e}}1 {í}{{\'i}}1 {ó}{{\'o}}1 {ú}{{\'u}}1
        {Á}{{\'A}}1 {É}{{\'E}}1 {Í}{{\'I}}1 {Ó}{{\'O}}1 {Ú}{{\'U}}1
        {ã}{{\~a}}1 {õ}{{\~o}}1 {Ã}{{\~A}}1 {Õ}{{\~O}}1
        {â}{{\^a}}1 {ê}{{\^e}}1 {ô}{{\^o}}1
        {à}{{\`a}}1 {ç}{{\c{c}}}1 {Ç}{{\c{C}}}1
}

\begin{document}

\section{Questão 08: Sistema de diagnóstico com redes bayesianas}

\subsection{Domínio e modelagem}

A rede estima a probabilidade de influenza em um cenário fictício. As cinco variáveis são binárias. E indica contato recente conhecido com pessoa infectada. I indica presença de influenza. F, T e M indicam febre, tosse e dor muscular. A seleção dos sintomas se baseia na descrição do CDC [1]. Os valores das tabelas foram definidos para fins educacionais.

\begin{figure}[htbp]
    \centering
    \begin{tikzpicture}[
        nodebox/.style={draw, rounded corners=2pt, minimum width=3.8cm, minimum height=0.65cm, align=center},
        >=Stealth
    ]
        \node[nodebox] (E) at (0, 2.8) {Exposição (E)};
        \node[nodebox, fill=blue!8] (I) at (0, 1.4) {Influenza (I)};
        \node[nodebox, minimum width=3.4cm] (F) at (-4.3, 0) {Febre (F)};
        \node[nodebox, minimum width=3.4cm] (T) at (0, 0) {Tosse (T)};
        \node[nodebox, minimum width=3.8cm] (M) at (4.5, 0) {Dor muscular (M)};

        \draw[->] (E) -- (I);
        \draw[->] (I) -- (F);
        \draw[->] (I) -- (T);
        \draw[->] (I) -- (M);
    \end{tikzpicture}
\end{figure}

A exposição aumenta a chance de influenza no modelo, e a doença influencia os três sintomas. A rede supõe que, conhecido o estado de I, os sintomas são mutuamente independentes e independentes de E. Sem condicionar em I, os sintomas podem estar associados pela causa comum. A distribuição conjunta se fatoriza da seguinte forma.

\[
P(E, I, F, T, M) = P(E)\,P(I \mid E)\,P(F \mid I)\,P(T \mid I)\,P(M \mid I)
\]

\subsection{Tabelas de probabilidades condicionais (CPTs)}

\begin{table}[htbp]
    \centering
    \begin{tabular}{>{\raggedright\arraybackslash}p{3.4cm} >{\raggedright\arraybackslash}p{4.6cm} >{\raggedright\arraybackslash}p{2.2cm} >{\raggedright\arraybackslash}p{2.2cm}}
        \toprule
        \textbf{Variável} & \textbf{Condição} & \textbf{P(Sim)} & \textbf{P(Não)} \\
        \midrule
        Exposição & Sem pais & 0,20 & 0,80 \\
        \addlinespace
        Influenza & Exposição = Não & 0,05 & 0,95 \\
        \addlinespace
        Influenza & Exposição = Sim & 0,30 & 0,70 \\
        \addlinespace
        Febre & Influenza = Não & 0,20 & 0,80 \\
        \addlinespace
        Febre & Influenza = Sim & 0,75 & 0,25 \\
        \addlinespace
        Tosse & Influenza = Não & 0,35 & 0,65 \\
        \addlinespace
        Tosse & Influenza = Sim & 0,80 & 0,20 \\
        \addlinespace
        Dor muscular & Influenza = Não & 0,15 & 0,85 \\
        \addlinespace
        Dor muscular & Influenza = Sim & 0,65 & 0,35 \\
        \bottomrule
    \end{tabular}
\end{table}

A exposição é tratada como menos frequente que sua ausência e eleva a probabilidade da doença de 5\% para 30\%. A tosse foi modelada como o sintoma mais comum nos dois grupos. Febre e dor muscular têm menor frequência de fundo. As probabilidades ficam entre 0 e 1 para admitir sintomas por outras causas e quadros de influenza sem todos os sintomas.

\subsection{Inferência exata}

O programa lê o grafo e as CPTs do arquivo \texttt{modelo\_influenza.json}. Ele enumera as 32 configurações binárias e calcula a massa conjunta de cada uma. Uma evidência observada elimina configurações incompatíveis. A massa das configurações restantes com I = Sim é dividida pela massa total restante. Variáveis sem observação são somadas sobre os dois estados.

Sem observações, $P(\mathrm{I} = \mathrm{Sim}) = 0,20 \times 0,30 + 0,80 \times 0,05 = 0,10$. Com febre observada, $P(\mathrm{I} = \mathrm{Sim} \mid \mathrm{F} = \mathrm{Sim}) = (0,10 \times 0,75) / (0,10 \times 0,75 + 0,90 \times 0,20) = 29,41\%$. A massa conjunta sem evidências foi verificada e soma 1.

Um relato incerto pode atribuir pesos de verossimilhança a ambos os estados. Para dor muscular, os pesos 0,2 (Não) e 0,8 (Sim) tornam o relato quatro vezes mais provável se o sintoma estiver presente. Cada configuração é multiplicada pelo peso correspondente antes da normalização. Os pesos descrevem a compatibilidade do relato com o estado e não a probabilidade posterior do sintoma.

\subsection{Experimentos}

Todos os cenários começam com a probabilidade a priori de 10,00\%. As evidências de uma mesma linha são acumuladas na ordem apresentada.

\begin{table}[htbp]
    \centering
    \begin{tabular}{>{\raggedright\arraybackslash}p{4.8cm} >{\raggedright\arraybackslash}p{10.0cm}}
        \toprule
        \textbf{Cenário} & \textbf{Evidências e atualizações de P(I = Sim)} \\
        \midrule
        1. Evidência incompleta & Febre = Sim: 29,41\%. Tosse e dor muscular ficam desconhecidas. \\
        \addlinespace
        2. Sintomas sucessivos & Febre = Sim: 29,41\%. Acrescenta tosse = Sim: 48,78\%. Acrescenta dor muscular = Sim: 80,50\%. \\
        \addlinespace
        3. Exposição e ausência de febre & Exposição = Sim: 30,00\%. Acrescenta febre = Não: 11,81\%. Acrescenta tosse = Sim: 23,44\%. \\
        \addlinespace
        4. Relato incerto & Febre = Sim: 29,41\%. Acrescenta relato de dor muscular com pesos 0,2 e 0,8: 45,88\%. \\
        \bottomrule
    \end{tabular}
\end{table}

A dor muscular observada teve o maior efeito na segunda sequência, pois é proporcionalmente mais frequente com influenza que sem a doença. No cenário 3, a ausência observada de febre reduz o resultado, e a tosse volta a elevá-lo. No cenário 4, o relato incerto altera a estimativa, mas tem efeito menor que uma observação positiva confirmada.

\subsection{Discussão e demonstração}

Em comparação com regras determinísticas, a rede combina sinais conflitantes, distingue ausência de falta de informação e produz uma estimativa gradual. As CPTs são hipotéticas, e a independência entre sintomas pode superestimar a força de evidências correlacionadas. Uma extensão poderia incluir outras doenças respiratórias, vacinação e resultados de exames, com parâmetros estimados e calibrados em dados adequados. O protótipo tem finalidade acadêmica.

Execute \texttt{python rede\_bayesiana.py} para reproduzir os quatro cenários. Para uma consulta própria, use, por exemplo, \texttt{python rede\_bayesiana.py febre=sim tosse=nao}.

\section*{Referências}

\noindent [1] CDC. \textbf{Signs and Symptoms of Flu}. Disponível em: \url{https://www.cdc.gov/flu/signs-symptoms/index.html}.

\end{document}
